\chapter{Referencial teórico}
\label{chap:referencial}

Este capítulo apresenta o conjunto de conceitos fundamentais necessários para
compreender a análise conduzida nesta dissertação. O capítulo discute: (i) a
formação da volatilidade implícita em mercados de opções; (ii) a mensuração da
volatilidade histórica; (iii) a definição, as interpretações e as propriedades
do prêmio de risco de volatilidade (VRP) e do prêmio do Bitcoin (BVRP),
incluindo os regimes de volatilidade; (iv) peculiaridades do mercado de opções
de Bitcoin; e (v) abordagens de aprendizado de máquina --- incluindo modelos
lineares regularizados (LASSO e Ridge) e métodos baseados em árvores (floresta
aleatória e XGBoost, implementação de \textit{gradient boosting}) --- aplicadas à previsão do BVRP e de
retornos do Bitcoin.

\section{Volatilidade implícita}

A volatilidade implícita (\textit{Implied Volatility}, IV) representa a
expectativa de mercado, no mundo neutro ao risco, sobre a variabilidade futura
do preço do ativo subjacente: $\text{IV}_t = E^*(\text{RV}_{t+1:t+h} \mid I_t)$,\footnote{A
rigor, o DVOL, como o VIX, é a raiz da variância esperada no mundo neutro ao
risco, $\text{IV}_t = \sqrt{E^*(\text{RV}^2_{t+1:t+h} \mid I_t)}$; a notação do
texto ignora a diferença de Jensen entre a raiz da esperança e a esperança da
raiz.}
em que o asterisco denota a expectativa no mundo neutro ao risco e $I_t$, a
informação disponível em $t$ (Seção~\ref{sec:refer-vrp-teorico}). Nesta dissertação,
a IV não vem da inversão de um modelo de precificação, como o de Black--Scholes:
o índice DVOL da Deribit, assim como o VIX, é uma estimativa livre de modelo,
extraída diretamente dos preços de opções fora do dinheiro
\citep{carr2009variance, cboe2019vix} (Seção~\ref{sec:refer-dvol}).

Diferentemente da volatilidade histórica, a volatilidade implícita não é uma
estimativa estatística no sentido clássico. Ela pode ser interpretada como um
\textit{preço} de risco associado à distribuição neutra ao risco dos retornos.
Como a estimativa livre de modelo mede a variação quadrática esperada dos
retornos, e não apenas a volatilidade de um processo de difusão contínuo, ela
incorpora também os saltos e, por isso, reflete assimetria, curtose e risco de
cauda \citep{carr2009variance}. Assim, fatores como os seguintes influenciam a IV:

\begin{itemize}
    \item \textbf{demanda por proteção}: a procura por proteção contra quedas,
          feita com opções de venda ou, pela paridade de compra e venda, com
          opções de compra \citep{hull2018options}, entra na formação dos
          preços das opções e, portanto, da IV, sobretudo em preços de
          exercício fora do dinheiro;
    \item \textbf{estrutura de \textit{skew}/\textit{smile}}: a IV varia entre
          preços de exercício e vencimentos. Parte dessa variação vem da
          própria distribuição do subjacente: quando a IV é calculada por
          Black--Scholes, que supõe retornos sem saltos, a assimetria e os
          saltos negativos exigem IV maiores nos preços de exercício baixos,
          o que gera o sorriso de volatilidade \citep{pan2006jumpRisk,
          gatheral2011volsurface};
    \item \textbf{eventos esperados}: anúncios regulatórios, atualizações de rede,
          vencimentos relevantes e mudanças de regime percebidas;
    \item \textbf{condições de liquidez}: diferença entre os preços de compra e
          de venda (\textit{bid--ask spread}), profundidade
          do livro de ofertas e variações intradiárias de fluxo.
\end{itemize}

Em mercados tradicionais, a superfície de volatilidade apresenta padrões
estilizados bem documentados, como \textit{smile} e \textit{skew}, associados a
distribuições assimétricas e a risco de cauda. No caso do Bitcoin, a curva de
volatilidade implícita tende a apresentar inclinação mais acentuada e maior
sensibilidade a fluxos de curto prazo, refletindo a elevada incerteza do ativo,
a negociação contínua e a menor maturidade relativa do mercado, em linha com a
literatura sobre superfícies de volatilidade \citep{gatheral2011volsurface}.

\section{Volatilidade histórica}

A volatilidade histórica (VH) é uma medida \textit{ex post} da variabilidade
dos retornos observados de um ativo ao longo do tempo, construída sob a medida
física. Trata-se de um estimador não paramétrico da volatilidade, amplamente
utilizado na literatura empírica, e constitui um insumo fundamental para a
comparação com a volatilidade implícita.

Esta dissertação calcula a volatilidade histórica a partir de retornos
logarítmicos diários, numa janela móvel de \(h\) dias:

\[
\text{VH}_t^{(h)} = \sqrt{ \frac{365}{h} \sum_{j=0}^{h-1} r_{t-j}^2 },
\]
em que \(r_{t-j} = \ln\left(\frac{P_{t-j}}{P_{t-j-1}}\right)\) representa o retorno
logarítmico diário do ativo, e o fator de anualização por 365 dias reflete a
negociação contínua do mercado de Bitcoin.

O nome segue a distinção usual da literatura. A volatilidade realizada
(\textit{realized volatility}, RV), no sentido de \citet{andersen2003realized},
estima a volatilidade de um período somando os retornos quadráticos de
frequência mais alta observados dentro dele --- por exemplo, retornos
intradiários para a volatilidade de um dia. A premissa é que a volatilidade
muda ao longo do tempo, e o estimador mede a volatilidade integrada no
intervalo. O estimador de volatilidade histórica, por sua vez, trata a
volatilidade como constante dentro da janela e tira a média dos retornos
quadráticos. A VH é, assim, o estimador com dados diários da RV que entra na
definição teórica do prêmio (Seção~\ref{sec:refer-vrp-teorico}); por isso, as duas
medidas aparecem com símbolos distintos.

A escolha de \(h\) muda as propriedades da medida. Com \(h\) grande, a janela
troca apenas uma de \(h\) observações de um dia para o outro, e a série fica
suave e reflete a volatilidade de um período longo; com \(h\) pequeno, a medida
é mais local e mais ruidosa. Esta dissertação usa \(h = 30\) dias corridos no
prêmio, o mesmo horizonte da volatilidade implícita com que a VH é comparada;
janelas de 1, 5, 60 e 90 dias aparecem apenas como variáveis explicativas nos
capítulos de previsão.

Alternativamente, pode-se definir a variância histórica como:

\[
\left(\text{VH}_t^{(h)}\right)^2 = \frac{365}{h} \sum_{j=0}^{h-1} r_{t-j}^2,
\]
sendo a volatilidade histórica a raiz quadrada dessa quantidade.

Embora a volatilidade histórica seja usualmente construída a partir de retornos
logarítmicos simples, pode-se alternativamente considerar retornos em excesso
em relação à taxa livre de risco, definidos como:

\[
r_t^e = r_t - r_{f,t},
\]
em que \(r_{f,t}\) representa a taxa livre de risco no período \(t\).

Nesse caso, a volatilidade histórica pode ser definida como:

\[
\text{VH}_t^{(h)} = \sqrt{ \frac{365}{h} \sum_{j=0}^{h-1} \left( r_{t-j} - r_{f,t-j} \right)^2 }.
\]

No entanto, para frequências diárias e, em particular, no contexto de criptoativos,
a diferença empírica entre utilizar retornos simples ou em excesso tende a ser
negligenciável. Por essa razão, a literatura usualmente adota retornos não ajustados,
abordagem também seguida neste trabalho.

A volatilidade realizada propriamente dita exigiria dados intradiários, o que
aumentaria a precisão (\textit{high-frequency realized volatility}); esta
dissertação usa a volatilidade histórica com dados diários, consistente
com a disponibilidade de dados do mercado à vista e com as séries derivadas do
ecossistema da Deribit, incluindo o índice DVOL.

A literatura clássica de volatilidade realizada, especialmente os trabalhos de
\citet{andersen2003realized} e \citet{corsi2009har}, estabelece a base teórica
para a construção de estimadores de volatilidade, além de documentar a
presença de persistência temporal e estrutura de memória longa na dinâmica da
volatilidade, fatos estilizados também observados em criptoativos.
\section{Prêmio de risco de volatilidade}
\label{sec:refer-vrp-teorico}

O prêmio de risco de volatilidade (VRP) pode ser interpretado como a
compensação exigida pelos investidores para assumir risco de volatilidade.
Formalmente, um prêmio de risco é a diferença entre a expectativa no mundo
físico e a expectativa no mundo neutro ao risco. Para um horizonte de $h$ dias,

\begin{equation}
\text{VRP}_t = E(\text{RV}_{t+1:t+h} \mid I_t) - E^*(\text{RV}_{t+1:t+h} \mid I_t),
\label{eq:refer-vrp-def}
\end{equation}
em que $\text{RV}_{t+1:t+h}$ é a volatilidade realizada entre $t+1$ e $t+h$,
$I_t$ é a informação disponível em $t$ e o asterisco denota a expectativa no
mundo neutro ao risco. O segundo termo é a volatilidade implícita,
$\text{IV}_t = E^*(\text{RV}_{t+1:t+h} \mid I_t)$.

O BVRP é o prêmio de risco de volatilidade do Bitcoin, com $h = 30$ dias. Como
a expectativa física não é observável, a variável-alvo desta dissertação é a
realização do prêmio, medida pela volatilidade histórica dos 30 dias seguintes:

\[
\text{BVRP}_t = \text{VH}_{t+1:t+30} - \text{IV}_t,
\]
em que $\text{VH}_{t+1:t+30}$ é a volatilidade histórica calculada com os
retornos de $t+1$ a $t+30$ e $\text{IV}_t$ é a volatilidade implícita a 30 dias
observada em $t$, dada pelo índice DVOL da Deribit. A esperança condicional de
$\text{BVRP}_t$ em $t$ é o VRP da equação~\eqref{eq:refer-vrp-def}, a menos da
diferença entre a VH e a RV. O BVRP de $t$ só é conhecido em $t+30$; por isso,
nada que seja usado como informação disponível em $t$ nos capítulos seguintes
depende dele.

Uma alternativa comum na literatura é substituir a expectativa física pela
volatilidade observada no passado recente: \citet{bollerslev2009expected}, por
exemplo, medem o prêmio de variância com a variância realizada do mês
anterior. A hipótese subjacente é a de passeio aleatório sem deriva para a
volatilidade --- a melhor previsão da VH dos próximos 30 dias é a VH dos
últimos 30, $E(\text{VH}_{t+1:t+30} \mid I_t) = \text{VH}_{t-29:t}$. Sob essa
hipótese, o prêmio esperado é aproximado por

\[
\text{BVRP}^{\text{proxy}}_t = \text{VH}_{t-29:t} - \text{IV}_t,
\]
conhecido em $t$. A hipótese é uma simplificação, e \citet{bekaert2014vix}
discutem suas limitações, estimando a expectativa física com modelos de
previsão da variância. Nesta dissertação, a proxy aproxima a expectativa do
prêmio, e não sua realização: ela entra como variável explicativa disponível
em $t$, enquanto o Capítulo~\ref{chap:previsao_bvrp} prevê o próprio BVRP.

A literatura empírica sobre prêmios de risco de volatilidade tem documentado esse fato estilizado em múltiplos mercados e horizontes. \citet{carr2009variance} formalizam a decomposição do retorno de variância em prêmio de risco e variação esperada, fornecendo o arcabouço metodológico que define o VRP como diferença entre medidas neutra ao risco e física. Para o mercado de ações, \citet{bollerslev2009expected} documentam, no contexto do S\&P~500, que o VRP é sistematicamente negativo e exibe poder preditivo para retornos de mercado em horizontes intermediários --- achado replicado em diversos índices acionários internacionais \citep{bollerslev2014international}. \citet{bali2009volatilitypremium} mostram, no corte transversal de ações individuais, uma relação negativa e significante entre o retorno esperado e o \textit{spread} entre a volatilidade realizada e a implícita. Os dois trabalhos que estudam o prêmio de volatilidade do Bitcoin com dados da Deribit \citep{alexanderImeraj2019, almeida2024cryptoVRP} e um trabalho relacionado \citep{han2026price} são discutidos na introdução (Capítulo~\ref{chap:introducao}), junto com o que esta dissertação acrescenta a eles.

Esta dissertação adota a convenção $\text{VH} - \text{IV}$ seguindo essa literatura. Sob esta convenção, valores negativos do BVRP correspondem ao estado típico em que a volatilidade implícita supera a volatilidade que efetivamente se realiza --- interpretação consistente com um prêmio pago pelos compradores de opções como proteção contra choques de variância. Quanto mais negativo o BVRP, maior o prêmio de seguro embutido nos preços de opções. Episódios de BVRP positivo ($\text{VH} > \text{IV}$) representam situações em que a volatilidade dos 30 dias seguintes supera a antecipada pelo mercado, tipicamente associadas a choques exógenos, liquidações em cascata ou movimentos abruptos não precificados.

A literatura também documenta que a volatilidade exibe persistência temporal elevada, com memória longa \citep{corsi2009har}, característica que tem implicações relevantes para a inferência estatística em modelos preditivos. Como VH e IV compartilham esse componente persistente, a diferença $\text{VH} - \text{IV}$ tende a cancelar boa parte dele; o Capítulo~\ref{chap:dados} mede esse padrão na amostra. No BVRP há, além disso, uma fonte mecânica de autocorrelação: alvos de datas vizinhas compartilham até 29 dos 30 retornos da janela. Esse padrão justifica os procedimentos econométricos adotados nos capítulos seguintes: erros-padrão de Newey--West calibrados ao horizonte, intervalos de confiança via reamostragem em blocos e janelas expansivas para qualquer cálculo de limiares estatísticos. A reamostragem em blocos aproxima a distribuição do estimador em amostra finita, o que permite avaliar a inferência quando há o viés de \citet{stambaugh1999}: quando o regressor de uma regressão preditiva é persistente e suas inovações se correlacionam com as do retorno, o estimador de MQO do coeficiente é viesado em amostras finitas, e o erro-padrão convencional não capta esse viés.

Do ponto de vista teórico, a distinção entre medida neutra ao risco (associada a
preços de opções) e medida física (associada a retornos observados) implica que
o prêmio incorpora, além de expectativas de variância,
componentes relacionados a risco de cauda e a preferências por proteção. Em
mercados tradicionais, como o S\&P 500, a evidência empírica documenta que o VRP
apresenta sinal sistematicamente negativo, associado à demanda por proteção contra
quedas abruptas e contra choques de volatilidade
\citep{bollerslev2011VRP, pan2006jumpRisk, carr2009variance, bollerslev2009expected}.

No mercado de Bitcoin, entretanto, \citet{alexanderImeraj2019} documentam,
sobretudo no início do mercado, que o BVRP apresenta dinâmica mais instável que a documentada no mercado de ações. Entre os
fatores frequentemente associados a essa instabilidade, estão:

\begin{enumerate}
    \item volatilidade histórica estruturalmente mais elevada, com episódios
          intermitentes de grande magnitude;
    \item eventos idiossincráticos e choques de liquidez em criptoativos, com
          impacto rápido sobre a volatilidade implícita;
    \item estrutura de mercado ainda em consolidação, com participação
          institucional relativamente menor que em mercados tradicionais e maior
          relevância de agentes alavancados;
    \item características microestruturais específicas do mercado de derivativos
          de cripto --- operação 24/7, contratos com liquidação inversa e
          mecanismos de auto-deleveraging --- que diferenciam a formação de preços
          de opções na Deribit dos mercados tradicionais regulados.
\end{enumerate}

Esses fatores pesam sobretudo em episódios de volatilidade elevada, o que
motiva separar a amostra em regimes de volatilidade. A análise usa dois
regimes, de baixa e de alta volatilidade, definidos pela volatilidade histórica
de 30 dias, cuja distribuição é bimodal na primeira parte da amostra. A
dissertação estima o corte entre os regimes apenas com dados da primeira janela
de estimação,
sem informação posterior às datas classificadas
(Capítulo~\ref{chap:metodologia}); no modelo em árvore do
Capítulo~\ref{chap:retornos_nao_linear}, o próprio modelo escolhe os cortes.

A combinação destes fatores produz três características empíricas relevantes para
a análise conduzida nesta dissertação: (i) o BVRP do Bitcoin é negativo em média
mas exibe distribuição assimétrica, com cauda direita mais pesada que no mercado de ações;
(ii) o sinal do prêmio é instável em horizontes curtos, com episódios
significativos de inversão; e (iii) os limiares estatísticos do BVRP variam
substancialmente ao longo do tempo, motivando o uso de janelas expansivas para
qualquer aplicação preditiva, conforme detalhado no Capítulo~\ref{chap:previsao_bvrp} e no
Apêndice~\ref{app:estrategias}.
\section{Mercado de opções de Bitcoin}

A Deribit domina amplamente o mercado de opções de Bitcoin e oferece liquidez
superior à das demais plataformas. Suas principais características incluem:

\begin{itemize}
    \item \textbf{negociação 24/7}: ausência de janelas de fechamento reduz o
          acúmulo de eventos fora do pregão, mas aumenta a relevância de efeitos
          intradiários de liquidez;
    \item \textbf{liquidez concentrada}: maior profundidade em opções
          no dinheiro e vencimentos curtos (7--30 dias), o que influencia
          a formação de IV em horizontes próximos;
    \item \textbf{efeito de liquidações automáticas}: posições alavancadas podem
          gerar movimentos bruscos e amplificar a volatilidade;
    \item \textbf{índice DVOL}: medida padronizada de volatilidade implícita
          calculada a partir de opções da Deribit.
\end{itemize}

A formação da IV no Bitcoin é sensível a fatores microestruturais, tais como:

\begin{itemize}
    \item tamanho mínimo de lote e restrições operacionais;
    \item diferenças mais amplas entre os preços de compra e de venda em
          horários de menor profundidade;
    \item heterogeneidade de formadores de mercado por região e perfil de risco;
    \item eventos associados ao \textit{funding} de futuros perpétuos, que afetam
          posicionamento e demanda por proteção.
\end{itemize}

Três características microestruturais merecem destaque adicional pelo impacto que exercem sobre a formação de preços de opções na Deribit. Primeiro, os contratos de opções são liquidados em Bitcoin, não em dólares --- propriedade conhecida como \textit{inverse settlement} --- o que introduz convexidade adicional no \textit{payoff} efetivo do detentor de opção quando expresso em moeda fiat. Segundo, a Deribit emprega mecanismo de \textit{auto-deleveraging} (ADL) como instrumento de gestão de risco: em situações de movimentos extremos em que o fundo de garantia da bolsa é insuficiente para cobrir as posições inadimplentes, a bolsa desfaz automaticamente contratos de contrapartes lucrativas. Esse mecanismo difere fundamentalmente dos \textit{circuit breakers} de bolsas tradicionais reguladas e introduz risco de contraparte sistemático que pode se refletir em prêmios elevados de volatilidade implícita em momentos de estresse. Terceiro, o mercado opera ininterruptamente, sem fechamentos noturnos ou de fim de semana, o que elimina o ``gap'' temporal entre informação e preço típico de mercados tradicionais, mas também concentra liquidez em horários de sobreposição com mercados asiáticos e norte-americanos.

A microestrutura aqui descrita passou por transformação institucional relevante durante o período amostral desta dissertação. Em 10 de janeiro de 2024, a \textit{Securities and Exchange Commission} (SEC) dos Estados Unidos aprovou a listagem de fundos negociados em bolsa (ETFs) que detêm Bitcoin à vista como ativo subjacente, contemplando onze produtos simultaneamente --- entre os quais o iShares Bitcoin Trust (IBIT, BlackRock), o Fidelity Wise Origin Bitcoin Fund (FBTC) e o ARK 21Shares Bitcoin ETF (ARKB). A aprovação marcou o ingresso pleno de canais institucionais regulados ao mercado de Bitcoin nos Estados Unidos. A implicação para o mercado de opções da Deribit é indireta mas substantiva: o aumento de fluxo institucional ao mercado à vista tende a comprimir a volatilidade histórica estrutural e a alterar a composição da demanda por proteção via opções, com potencial impacto sobre o nível e a persistência do BVRP --- canal cujo efeito a dissertação não testa e registra entre as limitações (Seção~\ref{sec:concl-limitacoes}).

Essas características tornam o BVRP um objeto empírico mais rico, porém mais
volátil, do que em mercados tradicionais. Na leitura desta dissertação, a
liquidez, a assimetria dos fluxos e a microestrutura também podem afetar a
formação da curva de IV no Bitcoin, um mercado menos maduro e mais fragmentado que os
tradicionais; sobre superfícies de volatilidade em geral, ver
\citet{gatheral2011volsurface}.

\subsection{O índice DVOL e a comparação com o VIX}
\label{sec:refer-dvol}

O índice DVOL da Deribit constitui medida de volatilidade implícita análoga, em sua intenção, ao VIX da Cboe. Conceitualmente, ambos buscam sintetizar a expectativa de mercado sobre a variabilidade futura do ativo subjacente em horizonte de 30 dias, extraída dos preços de opções vivas por um método livre de modelo. Formalmente, o DVOL é construído via integração ponderada dos preços de opções fora do dinheiro com vencimentos próximos a 30 dias, de modo análogo à formulação proposta por \citet{carr2009variance} para o conceito de variância neutra ao risco:

\begin{equation}
\text{DVOL}_t^2 = \frac{2}{\tau} \cdot e^{r\tau} \int_0^{\infty} \frac{Q(K)}{K^2} \, dK,
\label{eq:refer-dvol}
\end{equation}
em que $\tau = 30/365$, $Q(K)$ é o preço de mercado da opção fora do dinheiro com preço de exercício $K$ (opção de venda para $K < F$ e opção de compra para $K > F$, em que $F$ é o preço a termo), e $r$ é a taxa de juros livre de risco. A integral é aproximada na prática por uma soma discreta sobre os preços de exercício negociados, com interpolação para tratar lacunas na superfície de preços.

Quatro diferenças metodológicas e institucionais distinguem o DVOL do VIX. Primeiro, o subjacente do DVOL é o Bitcoin (preço à vista agregado da Deribit), enquanto o VIX reflete o índice S\&P~500 --- ativos com naturezas estruturalmente distintas em termos de regulação e composição de fluxo. Segundo, o DVOL utiliza opções da própria Deribit, plataforma não regulada por autoridade financeira tradicional, ao passo que o VIX emprega opções listadas na Cboe sob supervisão da SEC. Terceiro, o DVOL é atualizado continuamente em frequência intradiária (24/7), refletindo a operação contínua do mercado de Bitcoin; o VIX, em contraste, opera segundo o horário de pregão da Cboe. Quarto, a profundidade de preços de exercício ativos é menor no DVOL do que no VIX, o que pode amplificar ruído de microestrutura em períodos de baixa liquidez. Essas diferenças justificam tratar o DVOL como medida de volatilidade implícita do Bitcoin sem assumir equivalência direta com o VIX em propriedades estatísticas ou comportamento sob estresse de mercado.

\section{Aprendizado de máquina em finanças}

Modelos de aprendizado de máquina têm se tornado ferramentas cada vez mais
comuns na pesquisa empírica em finanças, especialmente em aplicações de previsão
de retornos e volatilidade. Em contextos com múltiplas variáveis explicativas,
possíveis não linearidades e interações complexas, métodos flexíveis podem
superar especificações lineares tradicionais. Trabalhos como \citet{Gu2020EmpiricalML}
mostram que modelos baseados em árvores, redes neurais e técnicas de regularização
são capazes de capturar relações não lineares e de alta dimensão entre
variáveis explicativas e retornos de ativos.

Esta dissertação emprega cinco modelos supervisionados, organizados em duas
famílias: modelos lineares --- MQO, Ridge e LASSO, em que os dois
últimos introduzem regularização sobre a versão não penalizada ---; e métodos
baseados em árvores (floresta aleatória e XGBoost),
utilizados pela capacidade de capturar não linearidades e interações entre
variáveis sem especificá-las de antemão, e de medir a importância de cada
variável explicativa na previsão. Em todos os modelos com hiperparâmetros, a validação
cruzada embargada faz a escolha a cada reestimação: a janela de treino de
cada dobra termina 30 dias antes da janela de validação, de modo que nenhum
alvo de treino se sobreponha aos de validação (Capítulo~\ref{chap:previsao_bvrp}).

\subsection{Modelos lineares: MQO, Ridge e LASSO}

Técnicas de regularização linear introduzem penalidades sobre os coeficientes
estimados com o objetivo de reduzir o sobreajuste e melhorar a
generalização fora da amostra, especialmente em contextos com múltiplos
preditores potencialmente correlacionados \citep{tibshirani1996lasso}.

O LASSO (\textit{Least Absolute Shrinkage and Selection Operator}) minimiza a
soma dos quadrados dos resíduos sujeita a uma penalidade $\ell_1$:
\[
\hat{\beta}^{\text{LASSO}} = \arg\min_{\beta}
    \left\{ \sum_{t=1}^{T}(y_t - \mathbf{x}_t'\beta)^2
            + \lambda \sum_{j=1}^{p} |\beta_j| \right\},
\]
em que $\lambda \geq 0$ é o parâmetro de regularização. A penalidade $\ell_1$
induz esparsidade: leva exatamente a zero os coeficientes associados a
preditores irrelevantes, o que faz uma seleção automática de variáveis.

O Ridge utiliza penalidade $\ell_2$:
\[
\hat{\beta}^{\text{Ridge}} = \arg\min_{\beta}
    \left\{ \sum_{t=1}^{T}(y_t - \mathbf{x}_t'\beta)^2
            + \lambda \sum_{j=1}^{p} \beta_j^2 \right\}.
\]
Diferentemente do LASSO, o Ridge retém todos os preditores, mas encolhe seus
coeficientes em direção a zero, sendo especialmente útil na presença de
multicolinearidade. Nesta dissertação, ambos os modelos funcionam como
modelos lineares de referência, intermediários entre a regressão por MQO
irrestrita e as abordagens baseadas em árvores.

\subsection{Floresta aleatória}

A floresta aleatória é um método de \textit{bagging} que cria múltiplas árvores de
decisão \citep{breiman2001randomforest}, cada uma treinada em uma amostra
\textit{bootstrap} dos dados. Trata-se de um dos algoritmos mais consagrados da
literatura de aprendizado de máquina, amplamente utilizado em aplicações
financeiras devido à sua capacidade de capturar relações não lineares e
interações complexas entre variáveis.

Árvores agregadas em estrutura de \textit{bagging} apresentam vantagens como:

\begin{itemize}
    \item captura de não linearidades;
    \item menor sensibilidade a valores extremos e a ruído;
    \item bom desempenho mesmo quando relações lineares são fracas.
\end{itemize}

Além disso, a aleatorização de subconjuntos de variáveis em cada divisão reduz a
correlação entre árvores e tende a melhorar a generalização fora da amostra.

\subsection{\textit{Gradient boosting} e XGBoost}

O \textit{gradient boosting} é uma técnica de \textit{ensemble} que constrói
árvores de decisão sequencialmente, de modo que cada nova árvore busca corrigir
os erros residuais das árvores anteriores \citep{friedman2001gradientboosting}.
Diferentemente da floresta aleatória, que utiliza \textit{bagging} com árvores
independentes, o \textit{gradient boosting} realiza otimização sequencial sobre
uma função de perda, resultando em maior sensibilidade a padrões sutis e
relações não lineares.

O modelo usado nesta dissertação é o XGBoost \citep{chen2016xgboost}, uma
implementação de \textit{gradient boosting}. A validação cruzada embargada
escolhe os hiperparâmetros --- número de árvores, taxa de aprendizado,
profundidade e peso mínimo por folha --- a cada reestimação, como nos demais
modelos (Capítulo~\ref{chap:previsao_bvrp}).

Entre suas características, estão:

\begin{itemize}
    \item construção sequencial de árvores para correção iterativa dos resíduos;
    \item boa capacidade de capturar não linearidades e interações entre
          variáveis;
    \item sensibilidade a padrões sutis, superior à de métodos de agregação
          independente como a floresta aleatória.
\end{itemize}

A dissertação emprega os cinco modelos --- MQO, LASSO, Ridge, floresta
aleatória e XGBoost --- para avaliar se o
BVRP é previsível fora da amostra a partir de variáveis de volatilidade, retorno
e estado de mercado, comparando modelos lineares e não lineares contra
modelos de referência. A floresta aleatória é também o modelo único em árvore
do Capítulo~\ref{chap:retornos_nao_linear}, que gera o BVRP previsto usado para testar a
relação não linear entre o prêmio e os retornos futuros.
